\chapter{Experimental Evidence}
\label{chap:experimental}

% ── Introduction ──────────────────────────────────────────────────────────────
\section{Introduction}
\label{sec:exp-intro}

This chapter surveys the experimental signatures predicted by the QOU
framework and quantum torsion theory developed in
Chapters~\ref{chap:qou}--\ref{chap:interpretations}, and reviews existing
observational data that may be relevant.

% ── Predicted Signatures ─────────────────────────────────────────────────────
\section{Predicted Observational Signatures}
\label{sec:exp-signatures}

% TODO: derive concrete observable signatures from the torsion tensor,
%       e.g., anomalous phase shifts, polarization rotations, or
%       deviations from standard dispersion relations.

\begin{conjecture}[Torsion Phase Shift]
  \label{conj:phase-shift}
  An interferometric experiment sensitive to the quantum connection
  $\nabla$ should detect a phase shift
  \begin{equation}
    \Delta\phi_{\mathcal{T}}
    = \frac{e}{\hbar} \oint_\gamma \mathcal{T}^\nabla,
    \label{eq:phase-shift}
  \end{equation}
  where the integral is taken over a closed loop $\gamma$ in the
  experimental apparatus.
\end{conjecture}

% ── Existing Data ─────────────────────────────────────────────────────────────
\section{Review of Existing Experimental Data}
\label{sec:exp-existing}

% TODO: survey neutron interferometry, atom interferometry, and CMB
%       polarization data for hints of torsional effects.

% ── Proposed Experiments ─────────────────────────────────────────────────────
\section{Proposed Experiments}
\label{sec:exp-proposed}

% TODO: outline specific, feasible experimental protocols that could
%       confirm or falsify Conjecture~\ref{conj:phase-shift}.

% ── Statistical Analysis ─────────────────────────────────────────────────────
\section{Statistical Analysis Framework}
\label{sec:exp-stats}

% TODO: describe the statistical tests, error budgets, and
%       Bayesian model-comparison methods appropriate for the proposed
%       experiments.

% ── Summary ───────────────────────────────────────────────────────────────────
\section{Summary}
\label{sec:exp-summary}

% TODO: summarize the experimental landscape and identify the most
%       promising near-term tests, linking to speculations in
%       Chapter~\ref{chap:speculations}.
